\chapter{Sequential Collection and Stopping of Evidence}
\label{sec:p15-sequential-evidence}

A sequential design specifies when to continue collecting observations
and how to decide when collection stops. Its error probabilities concern
that entire procedure. Wald's original treatment develops this distinction
for two simple hypotheses, derives likelihood-ratio boundaries, and
separates exact relations from approximations that neglect the final
boundary crossing \cite[Part I, Sections 2--4]{Wald1945SequentialTests}.

\section{Information, continuation and terminal decisions}
\label{sec:p15-se-information}

\begin{hbdefinition}{def:p15-se-model}{Two simple laws and their observed likelihoods}
Let \(P_0,P_1\) be two specified probability laws for an observation
on a measurable space, with densities \(f_0,f_1\) against a common
\(\sigma\)-finite measure \(\nu\); choosing \(\nu=P_0+P_1\)
always supplies one. Under hypothesis \(H_i\), observations are independent
with common law \(P_i\). Their observed information is
\(\mathcal F_m=\sigma(X_1,\ldots,X_m)\), with trivial initial
information. Put
\[
 p_{im}=\prod_{j=1}^mf_i(X_j),\qquad L_m=p_{1m}/p_{0m},\qquad L_0=1.
\]
The ratio is zero when only the numerator vanishes, infinity when only
the denominator vanishes, and one when both vanish, as in the source.
The latter event is null under either specified hypothesis. At a first
zero or infinite ratio the test below stops, so it never multiplies
opposite singular ratios after stopping. When the two laws are equivalent,
the scores and sums
\[
 z_j=\log\frac{f_1(X_j)}{f_0(X_j)},\qquad Z_m=\sum_{j=1}^m z_j=\log L_m
\]
are finite almost surely. Every finite-score construction under a third
law assumes that both chosen densities are finite and positive almost
surely under that law. Its score moments require their own assumptions;
they are not inherited from the two specified laws.
\end{hbdefinition}

\begin{hbdefinition}{def:p15-se-procedure}{A decision rule on finite histories}
For every \(m\geq1\), partition the measurable space of histories of
length \(m\) into terminal regions \(R_m^0,R_m^1\) and continuation
region \(R_m\). After a history that has not already stopped, choose
decision zero, decision one, or another observation according to those
regions. Let \(T\) be the first terminal stage, with
\(\inf\varnothing=\infty\). Thus \(\{T\leq m\}\in\mathcal F_m\).
On \(\{T=m\}\), the decision \(D\in\{0,1\}\) is determined by
that history; its events are measurable at stopping. A fixed sample of
size \(N\) is the special case with no terminal regions before \(N\)
and no continuation region at \(N\).

Write \(Q_i=\{T<\infty,D=i\}\) for the two terminal events and
\(\alpha'=P_0(Q_1)\), \(\beta'=P_1(Q_0)\) for the actual errors.
Nontermination is a separate outcome unless proved null. Targets
\(\alpha,\beta\) are prescribed numbers, not names for the actual
errors of an arbitrarily chosen design.
An explicitly randomized design may add an independent seed, with the
same specified distribution under both hypotheses, to the initial
information. The seed then cancels from every likelihood ratio, and
the decision regions are measurable in seed and observation history.
\end{hbdefinition}

\begin{hbtheorem}{thm:p15-se-posterior}{proved=annex:p15.se.posterior}{The source's posterior-odds motivation}
Suppose exactly \(H_0,H_1\) are possible, with positive prior
probabilities \(g_0,g_1\), \(g_0+g_1=1\). The posterior odds after
a finite history equal \((g_1/g_0)L_m\). For prescribed posterior
correct-decision levels \(d_0,d_1\in(1/2,1)\), the rules accepting
\(H_i\) when its posterior probability reaches \(d_i\) are
\[
 L_m\geq \frac{g_0d_1}{g_1(1-d_1)},\qquad
 L_m\leq \frac{g_0(1-d_0)}{g_1d_0},
\]
respectively. They are mutually exclusive. These are posterior
conditions under this specified prior model; a frequentist error target
is a different calibration \cite[Section 3.1]{Wald1945SequentialTests}.
\end{hbtheorem}

\begin{hbdefinition}{def:p15-se-sprt}{Wald's upper and lower likelihood boundaries}
The original general rule chooses \(0<B<A<\infty\). At stage
\(m\), accept \(H_1\) if
\(L_m\geq A\), accept \(H_0\) if \(L_m\leq B\), and continue
when \(B<L_m<A\). Equality therefore terminates. In the finite-score
case put \(a=\log A\), \(b=\log B\); then
\[
 T=\inf\{m\geq1:Z_m\notin(b,a)\},\qquad
 Q_1=\{Z_T\geq a\},\quad Q_0=\{Z_T\leq b\}.
\]
The letters \(A,B\) retain the source's likelihood-ratio meaning;
\(a,b\) are their logarithms. A choice of a maximum horizon or a
terminal decision at that horizon changes this procedure.
Except where the general joint-density or termination statement says
otherwise, the following calibration and operating calculations use
the source's practical case \(0<B<1<A\), hence \(b<0<a\).
\end{hbdefinition}

\section{Termination and exact error control}
\label{sec:p15-se-error-control}

\begin{hbtheorem}{thm:p15-se-termination}{proved=annex:p15.se.termination}{Finite expected termination under the two distinct simple laws}
For any finite \(0<B<A\), if \(P_0\ne P_1\), put
\(\eta_0=\int\sqrt{f_0f_1}\,d\nu<1\).
For every \(m\geq1\),
\[
 P_0(T>m)\leq B^{-1/2}\eta_0^m,\qquad
 P_1(T>m)\leq A^{1/2}\eta_0^m.
\]
Thus the test terminates almost surely with finite expected size under
both laws, including one-sided supports. If the laws coincide, their
ratio is one almost surely: the procedure never terminates when
\(B<1<A\), and otherwise terminates at its first observation.

For a third law under which finite scores are i.i.d., a sufficient
condition for finite expected termination is \(P(z_1\ne0)>0\).
No nonzero mean or finite variance is needed for that assertion. If the
score is identically zero, termination fails when \(B<1<A\) and
otherwise occurs at the first observation. These conditions fill
the stopping argument cited to cumulative-sum theory in source Section 3.2.
\end{hbtheorem}

\begin{hbtheorem}{thm:p15-se-errors}{proved=annex:p15.se.errors}{The exact two coupled error inequalities}
Under the two distinct simple laws,
\[
       A\alpha'\leq1-\beta',\qquad
       \beta'\leq B(1-\alpha').
\]
Consequently \(\alpha'\leq1/A\) and \(\beta'\leq B\).
In the error plane the allowed pair lies in the closed quadrilateral
with vertices
\[
 (0,0),\quad(1/A,0),\quad
 \left(\frac{1-B}{A-B},\frac{B(A-1)}{A-B}\right),\quad(0,B).
\]
For prescribed \(\alpha,\beta\in(0,1)\), choosing
\(A\geq1/\alpha\) and \(B\leq\beta\) guarantees both exact
targets. The proof compares the densities on disjoint finite stopping
histories; it does not assume that an unbounded stopped likelihood keeps
its mean without a limiting argument
\cite[Section 3.2, equations (3.16)--(3.21)]{Wald1945SequentialTests}.
\end{hbtheorem}

\begin{hbtheorem}{thm:p15-se-dependent}{proved=annex:p15.se.dependent}{The joint-density version for dependent observations}
Suppose \(p_{im}\) are densities of the actual joint histories under
two specified process laws, with their correct consistent marginals.
Apply the same ratio rule, with any \(0<B<A<\infty\), to
\(p_{1m}/p_{0m}\). If it terminates
almost surely under both process laws, the two coupled inequalities
above still hold. Independence is needed for the product-score and
termination arguments, not for this stopped-history density comparison.
Replacing a dependent joint density by a product of its marginals is
not this construction \cite[p.~130]{Wald1945SequentialTests}.
\end{hbtheorem}

\begin{hbtheorem}{thm:p15-se-nominal}{proved=annex:p15.se.nominal}{What the usual nominal boundaries actually imply}
For \(\alpha,\beta\in(0,1)\) and \(\alpha+\beta<1\), choose
\[
 A=\frac{1-\beta}{\alpha},\qquad B=\frac{\beta}{1-\alpha}.
\]
The resulting actual errors obey
\[
 \alpha'\leq\frac{\alpha}{1-\beta},\qquad
 \beta'\leq\frac{\beta}{1-\alpha},\qquad
 \alpha'+\beta'\leq\alpha+\beta.
\]
Either individual error may exceed its target. For
\(P_0=\operatorname{Bernoulli}(11/100)\),
\(P_1=\operatorname{Bernoulli}(99/100)\), and
\(\alpha=\beta=1/10\), these boundaries are \(A=9,B=1/9\).
The test stops after one observation, with
\(\alpha'=11/100>1/10\) and \(\beta'=1/100\).
The source's practical approximation is therefore not an exact
simultaneous error guarantee
\cite[Section 3.3, equations (3.26)--(3.30)]{Wald1945SequentialTests}.
\end{hbtheorem}

\section{Expected sample size and overshoot}
\label{sec:p15-se-sample-size}

\begin{hbdefinition}{def:p15-se-overshoot}{The quantities omitted by the boundary approximation}
Under a law with finite scores put \(\mu=Ez_1\), when this
expectation exists, and \(\gamma=P(Q_1)\). On a lower exit define
\(r_-=b-Z_T\geq0\); on an upper exit define
\(r_+=Z_T-a\geq0\). For events of positive probability write
\(c_-=E(r_-\mid Q_0)\), \(c_+=E(r_+\mid Q_1)\).
Products with a zero-probability terminal event are interpreted as zero.
These are conditional mean overshoots, not necessarily a common bound
on every observation.

For the one-observation score law, define
\[
 \xi=\sup_{u\geq0}E(z_1-u\mid z_1\geq u),\qquad
 \xi'=\inf_{u\geq0}E(z_1+u\mid z_1\leq-u),
\]
over conditioning events of positive probability. Allow infinite
extrema; useful finite bounds require \(\xi<\infty\) and
\(\xi'>-\infty\). The prime records the source's negative signed
quantity, so \(\xi'\leq0\) \cite[Section 4.1]{Wald1945SequentialTests}.
If one conditioning class is empty, set its corresponding bound to
zero; that terminal direction is impossible for the score law.
\end{hbdefinition}

\begin{hbtheorem}{thm:p15-se-mean}{proved=annex:p15.se.mean}{The stopped-sum identity with its integrability premises}
For i.i.d. finite scores, if \(E|z_1|<\infty\) and \(ET<\infty\),
then \(Z_T\) is integrable and
\[
 \mu ET=EZ_T=(1-\gamma)(b-c_-)+\gamma(a+c_+).
\]
When \(\mu\ne0\), dividing gives the exact expected size. Under the
finite extrema in the preceding definition,
\[
 (1-\gamma)(b+\xi')+\gamma a
 \leq \mu ET
 \leq (1-\gamma)b+\gamma(a+\xi).
\]
Division reverses the inequalities when \(\mu<0\). Neglecting
overshoot gives the approximation
\(ET\approx\{(1-\gamma)b+\gamma a\}/\mu\); it is not valid
by substitution when \(\mu=0\).
\cite[Equations (4.4)--(4.14)]{Wald1945SequentialTests}
\end{hbtheorem}

\begin{hbtheorem}{thm:p15-se-zero-drift}{proved=annex:p15.se.zero-drift}{A zero-mean score needs a second-moment identity}
In \cref{thm:p15-se-mean}, if additionally \(\mu=0\) and
\(0<\sigma_z^2=Ez_1^2<\infty\),
then
\[
 \sigma_z^2ET=EZ_T^2
 =(1-\gamma)E((b-r_-)^2\mid Q_0)
       +\gamma E((a+r_+)^2\mid Q_1).
\]
The terminal second moment is finite under these hypotheses.
If overshoot vanishes, then
\(\gamma=-b/(a-b)\) and \(ET=-ab/\sigma_z^2\).
An identically zero score has no exit and is excluded. This gives the
missing division-free case of the sample-size calculation.
\end{hbtheorem}

\section{Operating characteristics and exact finite calculations}
\label{sec:p15-se-operating-characteristics}

\begin{hbdefinition}{def:p15-se-exponential-root}{An exponential root and its one-step excess bounds}
For a third i.i.d. finite-score law, set \(\varphi(t)=Ee^{tz_1}\)
where finite. Suppose \(h\ne0\), \(\varphi(h)=1\), and
\(T<\infty\) almost surely. For conditioning events of positive
probability define
\[
 \eta=\inf_{v>1}vE(e^{hz_1}\mid e^{hz_1}\leq1/v),\qquad
 \delta=\sup_{0<v<1}vE(e^{hz_1}\mid e^{hz_1}\geq1/v).
\]
For the following finite bounds assume \(0<\eta\leq1\leq\delta<\infty\).
These source quantities differ from the affinity \(\eta_0\).
\end{hbdefinition}

\begin{hbtheorem}{thm:p15-se-oc}{proved=annex:p15.se.oc}{Exact exponential identity and bounded operating characteristic}
In the preceding setting, \(Ee^{hZ_T}=1\). If \(h>0\), then
\[
 \frac{1-B^h}{\delta A^h-B^h}
 \leq\gamma\leq
 \frac{1-\eta B^h}{A^h-\eta B^h}.
\]
If \(h<0\), then
\[
 \frac{1-A^h}{\delta B^h-A^h}
 \leq1-\gamma\leq
 \frac{1-\eta A^h}{B^h-\eta A^h}.
\]
The approximation obtained by setting the two conditional exponential
overshoots to their boundary values is
\[
 \bar\gamma=\frac{1-B^h}{A^h-B^h}.
\]
Its algebraic extension at \(h=0\) is \(-b/(a-b)\); this is an
overshoot approximation, not a claim that a nonzero root always exists.
\cite[Section 3.4, equations (3.32)--(3.44)]{Wald1945SequentialTests}
\end{hbtheorem}

\begin{hbtheorem}{thm:p15-se-lattice}{proved=annex:p15.se.lattice}{The source's finite lattice exit equations}
For positive integers \(g_1,g_2\), suppose the scores are integers
between \(-g_1\) and \(g_2\),
with both endpoint probabilities positive. Let \(q_j=P(z_1=j)\),
\(a_0=\lceil a\rceil\), \(b_0=\lfloor b\rfloor\).
The terminal values belong to
\[
 \{b_0-g_1+1,\ldots,b_0\}\ \cup\
 \{a_0,\ldots,a_0+g_2-1\}=\{c_1,\ldots,c_g\},
 \qquad g=g_1+g_2.
\]
For the \(g\) complex roots \(u_i\) of
\(\sum_{j=-g_1}^{g_2}q_j u^j=1\), counted with multiplicity,
form \(\mathcal M_{ij}=u_i^{c_j}\). If \(\det\mathcal M\ne0\), the terminal
probabilities are \(P(Z_T=c_j)=\det\mathcal M_j/\det\mathcal M\), where
\(\mathcal M_j\) replaces column \(j\) by ones. Summing the entries with
\(c_j\geq a_0\) gives \(\gamma\). Repeated roots may make this
particular determinant singular; no division by zero is permitted.
\cite[Equations (3.45)--(3.48)]{Wald1945SequentialTests}
\end{hbtheorem}

\begin{hbtheorem}{thm:p15-se-finite-recursion}{proved=annex:p15.se.finite-recursion}{Exact absorption recursion, including singular root matrices}
In the lattice setting let \(J=\{b_0+1,\ldots,a_0-1\}\).
The substochastic matrix \(K_J\) has entries
\(K_J(x,y)=P(z_1=y-x)\), \(x,y\in J\). Let
\(v_+(x)=P(x+z_1\geq a_0)\),
\(v_-(x)=P(x+z_1\leq b_0)\). Then
start the score walk at \(x\in J\), write \(P_x,E_x\) for its
probabilities and expectations, and let \(\gamma_x=P_x(Q_1)\).
Its exact absorption quantities are
\[
 \gamma_x=[(I-K_J)^{-1}v_+](x),\qquad
 E_xT=[(I-K_J)^{-1}{\bf1}](x).
\]
The inverse exists, and the same matrix determines finite-horizon
survival, both terminal decisions and all moments of the stopping time.
In particular, \(P_x(T>n)=[K_J^n{\bf1}](x)\).
Thus computation does not depend on the root determinant being nonzero.
For a finite nonlattice score alphabet, exact prefix enumeration through
a finite horizon is still possible; its unfinished probability is an
explicit error bound on either eventual decision probability.
\end{hbtheorem}

\section{Binomial and normal score formulas}
\label{sec:p15-se-score-families}

\begin{hbtheorem}{thm:p15-se-binomial}{proved=annex:p15.se.binomial}{The binomial score and its excess constants}
For \(0<p_0<p_1<1\), let \(P_i=\operatorname{Bernoulli}(p_i)\).
Under a third Bernoulli law with parameter \(p\in(0,1)\), the
score takes values
\[
 z_+=\log(p_1/p_0)>0,\qquad
 z_-=\log((1-p_1)/(1-p_0))<0.
\]
Here \(\mu=pz_++(1-p)z_-\), \(\xi=z_+\), \(\xi'=z_-\).
When \(\mu\ne0\), there is exactly one nonzero real exponential
root, opposite in sign to \(\mu\), determined by
\(pe^{hz_+}+(1-p)e^{hz_-}=1\). For \(h>0\),
\(\delta=e^{hz_+}\), \(\eta=e^{hz_-}\); for \(h<0\),
the two score values exchange roles. When \(\mu=0\), strict
convexity leaves only the root zero; the nonzero-root formula is not used.
\cite[Section 3.5 and Section 4.2]{Wald1945SequentialTests}
\end{hbtheorem}

\begin{hbtheorem}{thm:p15-se-normal}{proved=annex:p15.se.normal}{Normal scores, excess bounds and the zero-drift midpoint}
For unit-variance hypotheses with means \(-d,d\), \(d>0\), and
third law \(X\sim N(\theta,1)\), the score is \(z=2dX\).
Its mean is \(2d\theta\), variance \(4d^2\), and
\(\varphi(t)=\exp(2d\theta t+2d^2t^2)\).
For \(\theta\ne0\) its nonzero root is \(h=-\theta/d\).
Write \(\phi\) and \(\Phi\) for the standard normal density and
distribution function. Then
\[
 \xi=2d\left(\theta+\frac{\phi(\theta)}{\Phi(\theta)}\right),\qquad
 \xi'=-2d\left(\frac{\phi(\theta)}{1-\Phi(\theta)}-\theta\right).
\]
For the exponential bounds at \(\theta\ne0\),
\(\delta=\Phi(|\theta|)/(1-\Phi(|\theta|))\) and
\(\eta=1/\delta\). The annex proves the normal tail-ratio monotonicity
needed for these extrema. At \(\theta=0\) the mean is zero, and
the second-moment identity applies instead of division by the mean.
\cite[Equations (3.53)--(3.77), (4.17)--(4.24)]{Wald1945SequentialTests}
\end{hbtheorem}

\begin{hbmethod}{def:p15-se-normal-saving}{The source's approximate comparison with a fixed normal sample}
Assume \(\alpha,\beta\in(0,1)\) and \(\alpha+\beta<1\).
For known variance one and means
\(\theta_0<\theta_1\), let
\(q_u=\Phi^{-1}(u)\). The smallest fixed sample allowing errors no
larger than \(\alpha,\beta\) has size
\[
 N_{\rm fixed}=\left\lceil
 \frac{(q_{1-\alpha}+q_{1-\beta})^2}{(\theta_1-\theta_0)^2}
 \right\rceil .
\]
For the nominal SPRT thresholds, neglecting overshoot and replacing
actual errors by the targets gives
\[
 E_0T\approx
 \frac{-\{(1-\alpha)b+\alpha a\}}{(\theta_1-\theta_0)^2/2},\qquad
 E_1T\approx
 \frac{\beta b+(1-\beta)a}{(\theta_1-\theta_0)^2/2}.
\]
Without the integer rounding, the corresponding relative savings depend
only on the targets. At \(\alpha=\beta=.05\), the approximation
gives a saving of about \(51.03\%\) under either hypothesis, matching
the source's rounded Table 1 entry. It is not a uniform fifty-percent
guarantee. For separated means the omitted overshoot need not be small;
the exact inequalities and excess bounds above measure that omission.
\cite[Section 4.3, Table 1]{Wald1945SequentialTests}
\end{hbmethod}

\section{Stopping-time transforms and their approximation}
\label{sec:p15-se-time-transforms}

\begin{hbtheorem}{thm:p15-se-transform}{proved=annex:p15.se.transform}{The fundamental stopped transform with an explicit limiting condition}
Let finite scores be i.i.d., \(T<\infty\) almost surely, and
\(t\in\mathbb C\) satisfy \(Ee^{\operatorname{Re}(t)z}<\infty\),
\(\varphi(t)=Ee^{tz}\ne0\), \(|\varphi(t)|\geq1\).
Then, for bounded or unbounded scores,
\[
 E\left[e^{tZ_T}\varphi(t)^{-T}\right]=1.
\]
The one-step real-part exponential moment and finite expected exit
prove absolute integrability of the terminal integrand. Its surviving,
not-yet-stopped part is bounded by a constant and vanishes in expectation.
For \(t=0\) the identity is trivial; a nonzero real root reduces to
the proved exponential identity.
The same proof allows \(|\varphi(t)|<1\) if the terminal integrand
is absolutely integrable and its surviving expectation tends to zero;
a survival bound \(P(T>n)\leq C\rho^n\) with
\(\rho<|\varphi(t)|\) suffices for the latter condition.
\cite[Section 4.4, equation (4.35)]{Wald1945SequentialTests}
\end{hbtheorem}

\begin{hbmethod}{def:p15-se-transform-approximation}{The two-root boundary approximation}
Define \(\Psi_+(\tau)=E(e^{\tau T}\mid Q_1)\),
\(\Psi_-(\tau)=E(e^{\tau T}\mid Q_0)\) on their convergence
domains. The exponential transform is
\(\Psi(\tau)=\gamma\Psi_+(\tau)+(1-\gamma)\Psi_-(\tau)\).
Its restriction \(\Psi(is)\), \(s\in\mathbb R\), is the
characteristic function. Derivatives at real \(\tau=0\) give moments;
derivatives of \(\Psi(is)\) introduce the usual powers of \(i\).

Let \(t_1,t_2\) solve \(-\log\varphi(t)=\tau\), on selected local
branches tending to zero and the nonzero root. A consistent boundary
approximation replaces \(\gamma\) by \(\bar\gamma\) from
\cref{thm:p15-se-oc}, and replaces the conditional transforms by
\(\widehat\Psi_+,\widehat\Psi_-\). It uses the linear system
\[
 \bar\gamma A^{t_j}\widehat\Psi_+(\tau)
 +(1-\bar\gamma)B^{t_j}\widehat\Psi_-(\tau)=1\qquad(j=1,2).
\]
This calibration makes both conditional approximations equal one at
zero. The system is exact under zero overshoot when the stopped-transform
premises hold. For normally distributed scores with mean \(\mu\ne0\) and
variance \(\sigma_z^2>0\), the two roots are
\[
 t_{1,2}(\tau)=\frac{-\mu\mathbin\pm
                     \sqrt{\mu^2-2\sigma_z^2\tau}}{\sigma_z^2}.
\]
Set \(\widehat\Psi=\bar\gamma\widehat\Psi_+
 +(1-\bar\gamma)\widehat\Psi_-\).
Writing \(\mathcal D=A^{t_1}B^{t_2}-A^{t_2}B^{t_1}\), solving
the system gives this boundary approximation:
\[
 \widehat\Psi(\tau)=
 \frac{A^{t_1}+B^{t_2}-A^{t_2}-B^{t_1}}{\mathcal D}.
\]
At a vanishing denominator use a justified limit, not division by zero.
For \(\mu=0\), the two branches merge at zero; that case requires its
own limit and the second-moment calculation.

The source's one-sided continuous approximation, for positive drift
and no lower boundary, assigns to
\(m=\mu^2T/(2\sigma_z^2)\) the density
\[
 f_c(m)=\frac{c}{2\sqrt\pi\,m^{3/2}}
        \exp\left(-\frac{c^2}{4m}-m+c\right),\qquad
 c=\frac{\mu a}{\sigma_z^2}>0,\quad m>0.
\]
This is a continuous approximation to an integer stopping time, not its
exact probability mass function. The source's two-boundary series
construction is referred there to its separate 1944 cumulative-sum paper.
Small mean and standard deviation are the stated approximation regime;
they do not by themselves provide an explicit uniform error bound here.
\cite[Equations (4.36)--(4.52), pp.~148--151]{Wald1945SequentialTests}
\end{hbmethod}

\begin{hblemma}{thm:p15-se-continuous-density}{proved=annex:p15.se.continuous-density}{Normalization and moments of the continuous approximation}
For every \(c>0\), the function \(f_c\) in the preceding method is
a probability density. Its Laplace transform is
\[
       \int_0^\infty e^{-qm}f_c(m)\,dm
                   =\exp\{c(1-\sqrt{1+q})\}\qquad(q\geq0).
\]
It has mean \(c/2\) and variance \(c/4\). These are properties of
the continuous approximating density, not a proof of convergence to it
or a bound on its error for the original discrete stopping rule.
\end{hblemma}

\section{Finite horizons and repeated inspection}
\label{sec:p15-se-finite-horizons}

\begin{hbtheorem}{thm:p15-se-deadline}{proved=annex:p15.se.deadline}{What a fixed-horizon score says about earlier termination}
For the finite-score test and every \(N\geq1\),
\[
 P_0(T\leq N)\geq P_0(Z_N\leq b),\qquad
 P_1(T\leq N)\geq P_1(Z_N\geq a).
\]
If scores are normal with mean \(\mu_i\) and variance
\(\sigma_i^2>0\) under \(P_i\), these right sides are exactly
\[
 \Phi\left(\frac{b-N\mu_0}{\sigma_0\sqrt N}\right),\qquad
 1-\Phi\left(\frac{a-N\mu_1}{\sigma_1\sqrt N}\right).
\]
For other score laws a normal approximation requires additional
distributional justification. The source's Table 2 uses normal
hypotheses; it is not a table of the actual stopping-time distribution.
\cite[Section 4.5]{Wald1945SequentialTests}
\end{hbtheorem}

\begin{hbtheorem}{thm:p15-se-truncation}{proved=annex:p15.se.truncation}{A forced decision changes both error probabilities}
For the finite-score test, at a fixed horizon \(N\), keep every
earlier boundary decision. If
the original procedure would continue after observation \(N\), choose
\(D_N=1\) when \(Z_N>0\) and \(D_N=0\) when \(Z_N\leq0\).
Let \(\alpha_N,\beta_N\) be the errors of this truncated procedure.
Then
\[
 \alpha_N\leq\alpha'+P_0(0<Z_N<a),\qquad
 \beta_N\leq\beta'+P_1(b<Z_N\leq0).
\]
For normal scores the added terms are respectively
\[
 \Phi\left(\frac{a-N\mu_0}{\sigma_0\sqrt N}\right)
 -\Phi\left(\frac{-N\mu_0}{\sigma_0\sqrt N}\right),\qquad
 \Phi\left(\frac{-N\mu_1}{\sigma_1\sqrt N}\right)
 -\Phi\left(\frac{b-N\mu_1}{\sigma_1\sqrt N}\right).
\]
The base errors are the actual untruncated errors. Substituting nominal
targets is another approximation unless those base targets were separately
guaranteed. Safe upper substitutes are \(1/A\) and \(B\).
\cite[Section 4.6, equations (4.63)--(4.70)]{Wald1945SequentialTests}
\end{hbtheorem}

\begin{hbexample}{def:p15-se-normal-horizon}{Calculated termination and truncation quantities}
Take \(N(0,1)\) against \(N(\theta,1)\), with targets
\(\alpha=\beta=.05\) and
\(\theta=2\Phi^{-1}(.95)/\sqrt{1000}\approx.10402968\).
The fixed most-powerful design requires 1000 observations. Nominal
sequential boundaries have \(a=\log19\), \(b=-\log19\).
At horizon 1000 the preceding marginal calculation gives a termination
lower bound of about \(.773315\) under either law. The truncation
error addition is about \(.044456\) under either law. Adding a nominal
base \(.05\) produces the source-style approximation \(.094456\),
rounded upward to \(.095\) in Table 3. Using the guaranteed base
\(1/19\) instead gives the rigorous CDF-expression upper bound,
numerically about \(.097088\). These quantities are bounds or
approximations of the stated kind, not measured error frequencies.
\end{hbexample}

\begin{hbtheorem}{thm:p15-se-repeated-fixed}{proved=annex:p15.se.repeated-fixed}{Twenty valid fixed-sample inspections can produce sixty-four percent error}
Under a null model let \(U_1,\ldots,U_M\) be independent uniform
variables on \((0,1)\). The fixed one-observation rule rejects when
\(U_j\leq\alpha\), and has level exactly \(\alpha\) at each
preselected stage. The rule that rejects at the first such result through
stage \(M\) has type-I error
\[
                  1-(1-\alpha)^M.
\]
At \(\alpha=.05,M=20\) this is approximately \(.64151408\).
With unbounded inspection it is one. A sequentially valid comparison is
the specified first-crossing likelihood rule with \(A\geq1/\alpha\),
whose upper-exit error is at most \(\alpha\) over the entire horizon.
Independence is explicit in this calculated example; repeated cumulative
tests on overlapping samples have a different dependence structure and
require their own calculation.
\end{hbtheorem}

\section{The original efficiency result and its exact boundary}
\label{sec:p15-se-efficiency}

\begin{hbdefinition}{def:p15-se-efficiency}{Strength and expected-cost comparison}
In the source's terminology two terminating tests have equal strength
when their two actual errors agree. A test is stronger when neither
actual error is larger and at least one is smaller; crossed error
comparisons give no strength ordering. One is better if neither expected
size is larger and at least one is smaller. Admissibility excludes a
better test of equal strength. Minimizing the larger of the two expected
sizes is a different criterion from minimizing each coordinate.
For a test with finite worst expected size, its minimax efficiency is
the infimum of that worst expected size among equal-strength competitors,
divided by its own worst expected size. The definition uses an infimum
and does not assume that a minimizing competitor exists. Weighted
expected costs specify another criterion when the hypothesis weights
are supplied.
The original paper's general exact-optimality conjecture is not its
proved zero-overshoot result \cite[Sections 2.2 and 4.7]{Wald1945SequentialTests}.
\end{hbdefinition}

\begin{hbtheorem}{thm:p15-se-information}{proved=annex:p15.se.information}{The 1945 information lower bounds for any terminating test}
Suppose the two simple laws are equivalent and
\(E_i|z_1|<\infty\), with \(\mu_0=E_0z_1<0<\mu_1=E_1z_1\).
Any almost-surely terminating sequential procedure having actual errors
\(u,v\in(0,1)\) satisfies
\[
 E_0T\geq
 \frac{(1-u)\log((1-u)/v)+u\log(u/(1-v))}{-\mu_0},
\]
\[
 E_1T\geq
 \frac{v\log(v/(1-u))+(1-v)\log((1-v)/u)}{\mu_1}.
\]
Infinite expected sizes are allowed. The proof gives the finite-horizon
conditional-density argument and its truncation limit.
\cite[Section 4.7, equations (4.80)--(4.89)]{Wald1945SequentialTests}
\end{hbtheorem}

\begin{hbtheorem}{thm:p15-se-no-overshoot}{proved=annex:p15.se.no-overshoot}{Exact zero-overshoot efficiency and a quantified bounded-excess comparison}
Under the information-bound hypotheses, if this SPRT has no overshoot
at either boundary, it attains both lower bounds at its own actual
errors. It therefore minimizes both expected sizes among tests of
that same strength. This includes a nearest-neighbor score walk with
steps \(d,-d\) and boundaries at integral multiples of \(d\).

More generally suppose \(|z_1|\leq C<\infty\). Let \(I_i\) be
the relevant numerator in the preceding information bound, evaluated at
this test's actual errors. Then
\[
          0\leq E_iT-I_i/|\mu_i|\leq C/|\mu_i|.
\]
This bounds cost relative to the universal information lower bound.
The exact optimality conclusion above uses zero overshoot.
\end{hbtheorem}

The ordinary simple-law construction and the robust clipped-score tests
in \cref{sec:p15-robust-hypothesis-testing} have different model classes.
Their shared stopping identities do not make a nominal simple-law score
valid under an unspecified contamination or dependent process. For
unbounded entropy and singular robust endpoints, use the stated conditions
in \cref{sec:p15-robust-sequential-endpoints}; for the general stopping
identities, use \cref{sec:p15-optional-stopping-and-random-walks}.
\index{sequential testing}\index{optional stopping!error control}
\index{likelihood ratio!sequential boundaries}\index{overshoot}
